\scanpage{37}

$x\in U$, $\bar x^{(U)}\cap Y=\varnothing\ \Longrightarrow\
x\notin\operatorname{Ass}\underline O_X$, on conclut du

\uline{Lemme} Si $(Y,X)$ satisfait $(LG)_0$, alors pour
$\forall$ $x\in X$ tel que $Y\cap\bar x=\varnothing$,
on a $x\notin\operatorname{Ass}O_X$, du moins si $X$ est
quasi projectif.

\uline{Dém} Sinon, il existe un sous-faisceau coh $F$
\deleted{\unsure{de}} de $\underline O_X$
\deleted{\unsure{concentré}} de support $\bar x$.
Donc il existe une section non nulle de $F(n)\subset O_X(n)$,
[$n$ assez grand], \deleted{donc} \deleted{\unsure{qui}}
elle induira une section nulle de
$\underline O_{\deleted{\unsure{X}}}(n)_{/Y}$, absurde.

\fbox{\begin{minipage}{0.94\linewidth}
Conclusion pratique : Si $X$ quasi-projectif, alors
\diagram{\begin{tikzcd}[column sep=small,row sep=large,ampersand replacement=\&]
(LG)_0
\arrow[r,Leftarrow]
\arrow[r,Rightarrow,bend right=38,"\unsure{h}"'] \&
(LG)_{\text{faible}}
\arrow[r,Leftrightarrow]
\arrow[rr,Rightarrow,bend right=55,"\unsure{h}"'] \&
(LG)
\arrow[l,Rightarrow,bend left=38,"\unsure{h}"]
\arrow[r,Leftrightarrow] \&
\substack{H^0(X,\underline O_X(n))\xrightarrow{\sim}
H^0(X,\underline O_X(n)_{/Y})\\\text{pour }n\text{ grand}}
\end{tikzcd}}
\end{minipage}}

\margin{36}
